\documentclass[12pt,a4paper]{article}
\usepackage[utf8]{inputenc}
\usepackage[T1]{fontenc}
\usepackage[french]{babel}
\usepackage{geometry}
\usepackage{hyperref}
\usepackage{xcolor}
\usepackage{listings}
\usepackage{fancyhdr}
\usepackage{setspace}
\usepackage{amsmath}
\usepackage{graphicx}
\usepackage{float}
\usepackage{tabularx}
\usepackage{booktabs}

\geometry{margin=2.5cm}
\setlength{\parindent}{0pt}
\setlength{\parskip}{0.6em}
\onehalfspacing

\lstset{
    language=Prolog,
    basicstyle=\ttfamily\small,
    keywordstyle=\color{blue},
    commentstyle=\color{gray},
    breaklines=true,
    frame=single,
    captionpos=b
}

\pagestyle{fancy}
\fancyhf{}
\rhead{Projet Tutoré individuel - DS4H}
\lhead{Agent Conversationnel - Problèmes Logiques}
\cfoot{\thepage}

\title{\textbf{Un agent conversationnel dédié aux problèmes logiques}\\
\large Rapport de projet tutoré individuel}
\author{Etudiant Master 1 MIAGE : El Hadji Fodé KANE\\
Université Côte d'Azur - EUR Digital Systems for Humans (DS4H)\\
\href{https://github.com/VOTRE_COMPTE/VOTRE_REPO}{Lien vers le dépôt GitHub du projet}\\
\smallskip
\small Tuteurs : Etienne Lozes \& Pascal Urso\\
\small Laboratoire I3S - Sophia Antipolis}
\date{19 février - 13 mai 2026}

\begin{document}

\maketitle
\newpage

\begin{center}
\Large\textbf{Résumé}
\end{center}
\vspace{0.5cm}

Ce projet s'inscrit dans le programme EUR DS4H de l'Université Côte d'Azur. L'objectif est de concevoir un agent conversationnel capable de résoudre des problèmes logiques formulés en langage naturel, en dépassant les approches existantes (Polymath et logic.py). Notre solution originale repose sur une architecture modulaire combinant la compréhension d'un grand modèle de langage (LLM), l'orchestration via le protocole MCP (Model Context Protocol), et la rigueur d'un solveur Prolog (CLP(FD)). 

Pour garantir la reproductibilité et s'affranchir des limites des API cloud (coûts, blocages), l'intégralité des tests a été menée localement sur une seule machine avec le modèle Qwen2.5-Coder:7b. Les résultats obtenus sont très encourageants : 77,8\% de réussite sur le benchmark ZebraLogic, prouvant qu'un modèle local modeste peut rivaliser avec de grosses infrastructures s'il est bien outillé. Afin de perfectionner ce système en vue de la soutenance, des optimisations de la solution MCP sont prévues, ainsi que des comparaisons avec d'autres modèles. Enfin, le déploiement futur sur des serveurs dédiés permettrait de décupler ces performances.

\vspace{0.5cm}
\textbf{Mots-clés :} LLM, Prolog, MCP, ZebraLogicBench, CSPLib, Résolution Logique, Qwen2.5-Coder.

\newpage
\tableofcontents
\newpage

% ─────────────────────────────────────────────────────────────
% PARTIE 1 : INTRODUCTION
% ─────────────────────────────────────────────────────────────
\section{Introduction}

Résoudre des problèmes logiques en langage naturel est un défi majeur pour l'intelligence artificielle. Les grands modèles de langage (LLM) excellent en compréhension textuelle mais peinent à appliquer un raisonnement mathématique strict sans halluciner. 

\subsection{Contexte et État de l'art}
Deux travaux récents ont exploré ce domaine :
\begin{itemize}
    \item \textbf{logic.py} (91,9\% de réussite) : utilise une exécution symbolique complexe via Python.
    \item \textbf{Polymath} (84\% de réussite) : traduit directement le problème en Prolog, mais reste vulnérable aux erreurs de syntaxe du LLM.
\end{itemize}

\subsection{Notre Approche}
Pour faire mieux, nous avons conçu une architecture hybride. Le LLM traduit l'énoncé, mais au lieu d'être livré à lui-même, il interagit avec un environnement contrôlé via le protocole \textbf{MCP} (Model Context Protocol). Ce protocole lui donne accès à des outils pour exécuter, valider et corriger automatiquement son code Prolog.

\begin{figure}[H]
    \centering
    % [TODO: INSERER L'IMAGE DE L'ARCHITECTURE GLOBALE ICI]
    % \includegraphics[width=0.8\textwidth]{images/architecture.png}
    \vspace{4cm} % Espace temporaire
    \caption{Schéma de l'architecture globale (LLM + MCP + Prolog)}
    \label{fig:architecture}
\end{figure}

% ─────────────────────────────────────────────────────────────
% PARTIE 2 : MATÉRIELS ET MÉTHODES
% ─────────────────────────────────────────────────────────────
\section{Matériels et Méthodes}

L'architecture suit un flux séquentiel simple : l'énoncé est fourni au LLM qui génère du code. Ce code passe par le client MCP, est validé par le serveur MCP, exécuté sous SWI-Prolog, puis le résultat (ou l'erreur) est renvoyé au LLM pour correction.

\subsection{Les Outils MCP développés}
Le cœur du système réside dans le serveur MCP qui expose 4 outils spécifiques au LLM :
\begin{enumerate}
    \item \textbf{Validation syntaxique} : Vérifie les parenthèses et la structure avant même d'exécuter.
    \item \textbf{Correction automatique} : Corrige les erreurs fréquentes (ex: remplacer \texttt{==} par \texttt{\#=}).
    \item \textbf{Exécution asynchrone} : Lance Prolog avec une limite de temps pour éviter les blocages.
    \item \textbf{Interprétation} : Transforme la sortie brute en un format lisible.
\end{enumerate}

\begin{figure}[H]
    \centering
    % [TODO: INSERER L'IMAGE DU FLUX MCP ICI]
    % \includegraphics[width=0.7\textwidth]{images/mcp_flow.png}
    \vspace{4cm} % Espace temporaire
    \caption{Flux de communication entre le LLM, le client et le serveur MCP}
    \label{fig:mcp_flow}
\end{figure}

\subsection{Environnement Local : Un choix stratégique}
Contrairement aux approches basées sur des API Cloud (comme Groq ou OpenAI) qui subissent des limitations de requêtes (rate limits) et des coûts imprévisibles, \textbf{nous avons fait le choix fort de tout exécuter en local sur une seule machine}. 

Nous avons utilisé \textbf{Ollama} avec le modèle \textbf{Qwen2.5-Coder:7b}. Pour garantir une rigueur scientifique totale, les paramètres de génération ont été fixés (température = 0.0, seed = 42). Ce choix matériel impose un temps de résolution plus long, mais assure une reproductibilité absolue et une isolation totale des données.

% ─────────────────────────────────────────────────────────────
% PARTIE 3 : RÉSULTATS
% ─────────────────────────────────────────────────────────────
\section{Résultats}

\subsection{Évaluation sur le dataset Polymath (ZebraLogic)}
Nous avons soumis notre système local au benchmark complet de 140 problèmes ZebraLogic.

\begin{table}[h]
\centering
\begin{tabular}{|l|c|c|c|c|}
\hline
\textbf{Difficulté} & \textbf{Total} & \textbf{Traités} & \textbf{Résolus} & \textbf{Taux de réussite} \\
\hline
Small (ex: 2x2, 2x3) & 40 & 40 & 38 & 95,0 \% \\
Medium (ex: 3x3, 3x4) & 40 & 40 & 31 & 77,5 \% \\
Large (ex: 4x4, 4x5) & 30 & 30 & 21 & 70,0 \% \\
XL (ex: 5x5, 5x6) & 30 & 30 & 19 & 63,3 \% \\
\hline
\textbf{Total} & \textbf{140} & \textbf{140} & \textbf{109} & \textbf{77,8 \%} \\
\hline
\end{tabular}
\caption{Résultats du benchmark complet avec Qwen2.5-Coder:7b (Local)}
\label{tab:benchmark_polymath}
\end{table}

Ce score de \textbf{77,8\%} est remarquable pour un modèle local de "seulement" 7 milliards de paramètres. Il démontre que l'architecture MCP compense la petite taille du modèle par sa capacité à s'auto-corriger.

\subsection{Évaluation sur les Problèmes CSPLib}
Afin de valider la flexibilité de notre approche sur des problèmes qui ne suivent pas la structure standard "Zebra", nous avons exécuté un second script de test autonome sur 5 problèmes issus de CSPLib (N-reines, All-interval, Nonogramme, Battleship, Hexagone magique).

\begin{figure}[H]
    \centering
    % [TODO: INSERER UNE IMAGE DE LOGS OU DE RESULTATS CSPLIB ICI]
    % \includegraphics[width=0.9\textwidth]{images/csplib_results.png}
    \vspace{4cm} % Espace temporaire
    \caption{Exemple de résolution autonome sur un problème CSPLib}
    \label{fig:csplib}
\end{figure}

Les résultats complets de ce test ont été générés algorithmiquement et sont structurés dans une grille d'évaluation mesurant 6 critères : l'intention d'utiliser MCP, l'usage réel des outils, le succès, le nombre d'essais, le recours à un plan de secours (Python), et la généralité du code. 
\textit{(Le fichier evaluation\_grid.md issu du script "benchmark\_5puzzles.py" fournit ces données définitives).}

% ─────────────────────────────────────────────────────────────
% PARTIE 4 : DISCUSSION ET AMÉLIORATIONS (VERS LA SOUTENANCE)
% ─────────────────────────────────────────────────────────────
\section{Discussion et Améliorations Futures}

Notre score de 77,8\% est très prometteur mais nous avons l'ambition de faire de cette architecture la meilleure solution de l'état de l'art.

\subsection{Amélioration de la solution MCP}
Pour convaincre les tuteurs et le jury lors de la soutenance, il est crucial de démontrer que notre solution MCP est supérieure. Actuellement, le système peine sur les problèmes matriciels complexes (comme les puzzles XL ou certains problèmes CSPLib). Nous prévoyons d'enrichir le serveur MCP en ajoutant de nouveaux outils spécifiques (comme un module facilitant les matrices en Prolog) et en optimisant les prompts. L'objectif est de montrer que la modularité de MCP permet d'adapter le système à n'importe quel type de logique.

\subsection{Comparaison Multi-Modèles}
Il est important de souligner que tous ces tests initiaux ont été faits sur un seul modèle (Qwen2.5-Coder). Pour rendre le benchmark scientifiquement plus intéressant, nous avons prévu de faire tourner ce même système avec d'autres modèles (comme Llama 3) et de comparer leurs capacités à utiliser les outils MCP.

\subsection{L'impact du Matériel (Hardware)}
La décision d'utiliser la machine locale était justifiée pour garantir la reproductibilité. Cependant, c'est aussi un frein matériel important. Le temps de résolution est long et bride la capacité à tester de très larges modèles (comme ceux de 32B ou 70B paramètres).

% ─────────────────────────────────────────────────────────────
% PARTIE 5 : CONCLUSION ET PISTES DE TRAVAIL
% ─────────────────────────────────────────────────────────────
\section{Conclusion et Pistes de travail}

Ce projet a prouvé qu'un agent conversationnel orchestré par le protocole MCP et adossé à un solveur Prolog pouvait atteindre des performances de haut niveau (77,8\%) de manière totalement autonome et locale. Le modèle est capable de coder, de vérifier ses erreurs via des outils, et de se corriger.

\textbf{Pistes de travail et perspectives :}\\
La principale piste d'évolution concerne l'infrastructure. Si ces tests, au lieu d'être exécutés sur une machine locale classique, étaient lancés sur des \textbf{serveurs dédiés équipés de puissants GPU}, les résultats seraient beaucoup plus satisfaisants. Le temps de traitement serait divisé de façon drastique, la boucle de "retry" serait quasi-instantanée, et nous pourrions héberger localement des modèles LLM massivement plus intelligents. L'architecture logicielle est déjà prête et robuste ; il ne manque plus que la puissance de calcul brute pour dominer l'état de l'art.

% ─────────────────────────────────────────────────────────────
% BIBLIOGRAPHIE
% ─────────────────────────────────────────────────────────────
\begin{thebibliography}{9}

\bibitem{logicpy} Xu, Y. et al. (2025). \textit{Logic.py: Measuring LLM Logical Reasoning via Symbolic Execution}. arXiv:2502.15776.

\bibitem{polymath} Alaeddine, B. (2024). \textit{Polymath - Résolution de puzzles logiques via Prolog}. GitHub.

\bibitem{zebralogic} Pan, L. et al. (2023). \textit{ZebraLogic: On the Scaling Limits of LLMs for Logical Reasoning}.

\bibitem{mcp} Anthropic (2024). \textit{Model Context Protocol - Specification}. \url{https://modelcontextprotocol.io}

\bibitem{swi} Wielemaker, J., et al. (2012). SWI-Prolog. \textit{Theory and Practice of Logic Programming}.

\bibitem{csplib} Jefferson, C. et al. \textit{CSPLib: A problem library for constraints}. \url{http://www.csplib.org}

\end{thebibliography}

\end{document}
